\begin{longtable}{@{}p{0.17\textwidth}p{0.14\textwidth}p{0.54\textwidth}r@{\hspace{4pt}}r@{}}
\caption{Full hierarchical cognitive-bias inventory at leaf level}\label{tab:s10-taxonomy}\\
\toprule
Cluster & Cognitive bias & Definition & $n_c$ & $n_f$ \\
\midrule
\endfirsthead
\multicolumn{5}{@{}l}{\textbf{Table \thetable\ (continued)}}\\
\multicolumn{5}{@{}p{\textwidth}@{}}{Full hierarchical cognitive-bias inventory at leaf level}\\
\toprule
Cluster & Cognitive bias & Definition & $n_c$ & $n_f$ \\
\midrule
\endhead
\midrule
\multicolumn{5}{r}{\textit{Continued on next page}}\\
\endfoot
\bottomrule
\endlastfoot
\addlinespace[2pt]\multicolumn{5}{@{}l}{\textbf{Attention \& Salience} ($n_{\mathrm{leaves}}=20$)} \\
\noalign{\needspace{10\baselineskip}}
\multirow{8}{0.17\textwidth}{Attentional Selection Under Load} & Attentional Bias & Preferential processing of threat-related over neutral stimuli. & 8 & 20 \\
 & Attentional Narrowing & Focus contracts under stress, excluding peripheral information. & 8 & 20 \\
 & Change Blindness & Failure to detect visual changes across interruptions or edits. & 8 & 20 \\
 & Cognitive Tunneling & Fixation on one task element at the expense of others. & 8 & 20 \\
 & Divided Attention Effect & Performance degrades when resources split across concurrent tasks. & 8 & 20 \\
 & Goal Shielding & Active suppression of competing goals distorts relevant-signal detection. & 8 & 20 \\
 & Inattentional Blindness & Unexpected stimuli go unnoticed when attention is engaged elsewhere. & 8 & 20 \\
 & Vigilance Decrement & Sustained-attention accuracy declines over time on watch tasks. & 8 & 20 \\
\noalign{\needspace{8\baselineskip}}
\multirow{6}{0.17\textwidth}{Salience Prominence And Vividness} & Novelty Effect & New or unfamiliar stimuli capture disproportionate attentional weight. & 6 & 20 \\
 & Picture Superiority Effect & Images are encoded and recalled more easily than words. & 6 & 20 \\
 & Prominence Effect & More important attributes receive weight beyond their informational value. & 6 & 20 \\
 & Salience Bias & Perceptually conspicuous information dominates judgments regardless of relevance. & 6 & 20 \\
 & Vividness Effect & Concrete, emotionally charged information outweighs abstract statistical data. & 6 & 20 \\
 & Von Restorff Effect & Distinctive items in a set are better remembered than homogeneous ones. & 6 & 20 \\
\noalign{\needspace{8\baselineskip}}
\multirow{6}{0.17\textwidth}{Threat And Warning Signal Processing} & Alarm Fatigue & Repeated alerts reduce vigilance and response rates over time. & 6 & 20 \\
 & Banner Blindness & Habituation causes users to ignore interface warning banners. & 6 & 20 \\
 & Cry Wolf Effect & False alarms erode trust and reduce future warning compliance. & 6 & 20 \\
 & Negativity Bias & Negative information receives disproportionately greater attentional and decisional weight. & 6 & 20 \\
 & Threat Superiority Effect & Threat-relevant stimuli are detected faster than neutral equivalents. & 6 & 20 \\
 & Warning Habituation & Repetitive exposure to warnings reduces their perceived urgency. & 6 & 20 \\
\addlinespace[2pt]\multicolumn{5}{@{}l}{\textbf{Trust \& Credibility} ($n_{\mathrm{leaves}}=14$)} \\\nopagebreak
\noalign{\needspace{9\baselineskip}}
\multirow{7}{0.17\textwidth}{Familiarity Fluency And Repetition} & Familiarity Heuristic & Prior exposure increases perceived trustworthiness of a source or claim. & 7 & 14 \\
 & Fluency Heuristic & Easy-to-process information is judged as more credible or true. & 7 & 14 \\
 & Illusory Truth Effect & Repeated statements are rated as more truthful than novel ones. & 7 & 14 \\
 & Mere Exposure Effect & Repeated exposure increases positive evaluation without conscious recognition. & 7 & 14 \\
 & Processing Fluency Effect & High cognitive ease inflates confidence and credibility assessments. & 7 & 14 \\
 & Recognition Heuristic & Recognized options are judged superior even without additional knowledge. & 7 & 14 \\
 & Truth Bias & Default assumption that communications are honest and accurate. & 7 & 14 \\
\noalign{\needspace{4\baselineskip}}
\multirow{2}{0.17\textwidth}{First Impression And Surface Inference} & First Impression Bias & Initial encounters anchor subsequent evaluations of trustworthiness. & 2 & 14 \\
 & Primacy Effect & Information presented first disproportionately shapes overall impressions. & 2 & 14 \\
\noalign{\needspace{7\baselineskip}}
\multirow{5}{0.17\textwidth}{Source Credibility Shortcuts} & Expertise Heuristic & Apparent expertise signals substitute for content-based verification. & 5 & 14 \\
 & Halo Effect & Positive impressions in one domain generalize to unrelated domains. & 5 & 14 \\
 & Horn Effect & Negative impressions in one domain generalize negatively to others. & 5 & 14 \\
 & Prestige Bias & High-status sources are granted credibility beyond their actual track record. & 5 & 14 \\
 & Reputation Heuristic & Source reputation proxies for message accuracy without content scrutiny. & 5 & 14 \\
\addlinespace[2pt]\multicolumn{5}{@{}l}{\textbf{Evidence \& Belief} ($n_{\mathrm{leaves}}=38$)} \\
\noalign{\needspace{8\baselineskip}}
\multirow{6}{0.17\textwidth}{Anchoring And Framing} & Anchoring Effect & Initial numeric value unduly constrains subsequent estimates and adjustments. & 6 & 38 \\
 & Attribute Framing & Description of a single attribute as gain or loss alters choice. & 6 & 38 \\
 & Attribute Substitution & Difficult judgment replaced by easier, related attribute evaluation. & 6 & 38 \\
 & Framing Effect & Logically equivalent options differ in appeal based on presentation format. & 6 & 38 \\
 & Goal Framing & Behavior influenced by whether outcomes are framed as gains or losses. & 6 & 38 \\
 & Insufficient Adjustment & Starting-point anchor is insufficiently corrected despite motivation to adjust. & 6 & 38 \\
\pagebreak[4]
\multirow{8}{0.17\textwidth}{Causal Pattern And Story Completion} & Apophenia & Perception of meaningful patterns in random or unrelated data. & 8 & 38 \\
 & Causal Oversimplification & Complex causation reduced to a single convenient explanatory factor. & 8 & 38 \\
 & Clustering Illusion & Random sequences perceived as non-random due to local clustering. & 8 & 38 \\
 & Illusion Of Validity & Confidence in predictions exceeds what pattern consistency warrants. & 8 & 38 \\
 & Intentionality Bias & Events attributed to deliberate agency rather than chance or error. & 8 & 38 \\
 & Narrative Fallacy & Coherent stories constructed post-hoc to explain sequences of events. & 8 & 38 \\
 & Proportionality Bias & Large effects presumed to require proportionally large causes. & 8 & 38 \\
 & Teleological Bias & Tendency to view events as purposeful and goal-directed by nature. & 8 & 38 \\
\noalign{\needspace{10\baselineskip}}
\multirow{8}{0.17\textwidth}{Confirmatory Processing} & Attitude Polarization & Exposure to mixed evidence strengthens rather than moderates existing beliefs. & 8 & 38 \\
 & Belief Perseverance & Initial beliefs resist change even after the supporting evidence is discredited. & 8 & 38 \\
 & Biased Assimilation & Evidence consistent with prior beliefs is accepted more readily than contradicting evidence. & 8 & 38 \\
 & Confirmation Bias & Preferential search for and interpretation of information confirming existing beliefs. & 8 & 38 \\
 & Motivated Reasoning & Reasoning process is recruited to justify a desired conclusion. & 8 & 38 \\
 & Myside Bias & Arguments on one's own side evaluated more favorably than opposing arguments. & 8 & 38 \\
 & Selective Exposure & Seeking information environments that confirm pre-existing attitudes. & 8 & 38 \\
 & Selective Scrutiny & Counterattitudinal information subjected to more rigorous critical analysis. & 8 & 38 \\
\noalign{\needspace{11\baselineskip}}
\multirow{9}{0.17\textwidth}{Representa\-tiveness And Diagnosticity Errors} & Base Rate Neglect & Prior probability information discounted in favor of case-specific details. & 9 & 38 \\
 & Conjunction Fallacy & Conjunction of events judged more probable than its constituents. & 9 & 38 \\
 & Denominator Neglect & Numerator frequencies attended to while denominators are ignored. & 9 & 38 \\
 & Disjunction Fallacy & Probability of at least one event underestimated relative to its constituents. & 9 & 38 \\
 & Insensitivity To Sample Size & Small samples treated as equally diagnostic as large ones. & 9 & 38 \\
 & Law Of Small Numbers & Small samples expected to reflect population parameters as reliably as large ones. & 9 & 38 \\
 & Ratio Bias & Absolute frequencies preferred over equivalent ratios for risk judgments. & 9 & 38 \\
 & Regression Fallacy & Regression to the mean attributed to causal intervention rather than statistical chance. & 9 & 38 \\
 & Representativeness Heuristic & Typicality replaces probability as the primary basis for categorization judgments. & 9 & 38 \\
\noalign{\needspace{9\baselineskip}}
\multirow{7}{0.17\textwidth}{Updating Failures} & Belief Echo Effect & Corrected misinformation continues to influence reasoning after retraction. & 7 & 38 \\
 & Conservatism Bias & Posterior beliefs updated too slowly in response to new evidence. & 7 & 38 \\
 & Continued Influence Effect & Discredited information persists in shaping subsequent inferences. & 7 & 38 \\
 & Misinformation Effect & Post-event information corrupts the accuracy of original memory. & 7 & 38 \\
 & Order Effect & Sequential order of evidence presentation alters final judgment. & 7 & 38 \\
 & Premature Closure & Decision made before all relevant information has been considered. & 7 & 38 \\
 & Recency Bias & Most recently encountered information weighted disproportionately in judgment. & 7 & 38 \\
\addlinespace[2pt]\multicolumn{5}{@{}l}{\textbf{Memory \& Source} ($n_{\mathrm{leaves}}=14$)} \\\nopagebreak
\noalign{\needspace{5\baselineskip}}
\multirow{3}{0.17\textwidth}{Availability And Retrievability} & Availability Cascade & Repeated media coverage makes an event seem more probable. & 3 & 14 \\
 & Availability Heuristic & Ease of recalling examples used to judge frequency or probability. & 3 & 14 \\
 & Retrievability Bias & Memorable information over-weighted relative to less retrievable alternatives. & 3 & 14 \\
\noalign{\needspace{7\baselineskip}}
\multirow{5}{0.17\textwidth}{Reconstructive Memory Distortions} & Consistency Bias & Past attitudes recalled as more consistent with current views than they were. & 5 & 14 \\
 & Hindsight Bias & Outcomes seem predictable in retrospect after their occurrence is known. & 5 & 14 \\
 & Peak End Rule & Episode evaluated by its most intense and final moments only. & 5 & 14 \\
 & Rosy Retrospection & Past experiences recalled more positively than they were when experienced. & 5 & 14 \\
 & Telescoping Effect & Recent events perceived as more distant; distant events as more recent. & 5 & 14 \\
\noalign{\needspace{8\baselineskip}}
\multirow{6}{0.17\textwidth}{Source Monitoring And Attribution Errors} & Cryptomnesia & Forgotten memories recur as seemingly new, original ideas. & 6 & 14 \\
 & False Attribution & Memory content correctly recalled but incorrectly assigned to a source. & 6 & 14 \\
 & Misattribution & Memory confidently attributed to an incorrect origin or context. & 6 & 14 \\
 & Sleeper Effect & Discounted message source forgotten while persuasive content is retained. & 6 & 14 \\
 & Source Confusion & Failure to distinguish between actual and imagined or suggested sources. & 6 & 14 \\
 & Source Monitoring Error & Inability to correctly identify the origin of stored information. & 6 & 14 \\
\addlinespace[2pt]\multicolumn{5}{@{}l}{\textbf{Risk \& Probability} ($n_{\mathrm{leaves}}=24$)} \\
\noalign{\needspace{11\baselineskip}}
\multirow{9}{0.17\textwidth}{Certainty Ambiguity And Probability Weighting} & Ambiguity Aversion & Known risk preferred over equivalent unknown or ambiguous outcomes. & 9 & 24 \\
 & Certainty Effect & Certain outcomes weighted disproportionately relative to probable ones. & 9 & 24 \\
 & Denominator Insensitivity & Probability judgments insensitive to changes in the denominator. & 9 & 24 \\
 & Possibility Effect & Small but non-zero probabilities overweighted relative to zero. & 9 & 24 \\
 & Probability Neglect & Low-probability extreme outcomes ignored when affect is engaged. & 9 & 24 \\
 & Probability Weighting & Objective probabilities transformed nonlinearly into subjective decision weights. & 9 & 24 \\
 & Pseudo Certainty Effect & Conditional outcomes framed as certain produce the certainty effect. & 9 & 24 \\
 & Scope Insensitivity & Willingness to pay fails to scale proportionally with scope of benefit. & 9 & 24 \\
 & Zero Risk Bias & Preference for eliminating small risks entirely over larger proportional reductions. & 9 & 24 \\
\noalign{\needspace{5\baselineskip}}
\multirow{3}{0.17\textwidth}{Randomness Frequency And Small Sample Inference} & Base Rate Dilution & Irrelevant non-diagnostic information dilutes base rate utilization. & 3 & 24 \\
 & Gambler Fallacy & Independent random events believed to be influenced by prior outcomes. & 3 & 24 \\
 & Hot Hand Fallacy & Runs of success seen as predictive of continued future success. & 3 & 24 \\
\noalign{\needspace{8\baselineskip}}
\multirow{6}{0.17\textwidth}{Reference Dependence And Risk Attitude} & Endowment Effect & Owned items valued more than identical unowned items. & 6 & 24 \\
 & Loss Aversion & Losses weighted more heavily than equivalent-magnitude gains. & 6 & 24 \\
 & Omission Bias & Harmful inactions judged as less bad than equivalent harmful actions. & 6 & 24 \\
 & Reference Dependence & Outcomes evaluated relative to a reference point rather than absolutely. & 6 & 24 \\
 & Risk Aversion & Preference for certain outcomes over equivalent-value uncertain ones. & 6 & 24 \\
 & Status Quo Bias & Existing state preferred over change even when change is advantageous. & 6 & 24 \\
\noalign{\needspace{8\baselineskip}}
\multirow{6}{0.17\textwidth}{Vulnerability And Severity Estimation} & Dread Risk Effect & Low-probability catastrophic outcomes evoke disproportionate dread and avoidance. & 6 & 24 \\
 & Invulnerability Bias & Personal risk underestimated relative to risk for others in similar circumstances. & 6 & 24 \\
 & Normalcy Bias & Threat magnitude underestimated by assuming situations will remain normal. & 6 & 24 \\
 & Optimism Bias & Positive outcomes expected more likely for self than for others. & 6 & 24 \\
 & Pessimism Bias & Negative outcomes expected more likely than objective probabilities warrant. & 6 & 24 \\
 & Unrealistic Optimism & Future personal outcomes rated as better than statistical base rates predict. & 6 & 24 \\
\addlinespace[2pt]\multicolumn{5}{@{}l}{\textbf{Temporal \& Commitment} ($n_{\mathrm{leaves}}=17$)} \\
\noalign{\needspace{5\baselineskip}}
\multirow{3}{0.17\textwidth}{Control And Protective Action} & Action Bias & Doing something preferred over inaction even when inaction is optimal. & 3 & 17 \\
 & Illusion Of Control & Outcomes perceived as personally controllable when they are not. & 3 & 17 \\
 & Risk Compensation & Increased safety measures prompt riskier behavior to restore target risk level. & 3 & 17 \\
\noalign{\needspace{7\baselineskip}}
\multirow{5}{0.17\textwidth}{Escalation And Commitment} & Effort Justification & Greater effort invested makes outcomes seem more valuable post-hoc. & 5 & 17 \\
 & Entrapment & Continued resource investment driven by accumulated prior costs. & 5 & 17 \\
 & Escalation Of Commitment & Commitment increases to a failing course to justify prior decisions. & 5 & 17 \\
 & Plan Continuation Bias & Original plan pursued despite evidence that conditions have changed. & 5 & 17 \\
 & Sunk Cost Fallacy & Past irrecoverable costs drive future decisions rather than prospective value. & 5 & 17 \\
\noalign{\needspace{6\baselineskip}}
\multirow{4}{0.17\textwidth}{Inertia Delay And Nonaction} & Default Effect & Pre-selected option accepted without active deliberation. & 4 & 17 \\
 & Inaction Inertia & Missing one opportunity reduces the likelihood of taking a later one. & 4 & 17 \\
 & Inertia Bias & Preference for maintaining current behavior over switching alternatives. & 4 & 17 \\
 & Procrastination & Delay of intended actions despite awareness of negative consequences. & 4 & 17 \\
\noalign{\needspace{7\baselineskip}}
\multirow{5}{0.17\textwidth}{Present Oriented Choice} & Delay Discounting & Rewards lose value as a function of their temporal distance. & 5 & 17 \\
 & Hyperbolic Discounting & Discount rate declines steeply for near-future but shallowly for distant rewards. & 5 & 17 \\
 & Immediacy Effect & Present rewards weighted more than equivalent future rewards. & 5 & 17 \\
 & Myopic Loss Aversion & Short evaluation horizons amplify loss aversion and reduce risk-taking. & 5 & 17 \\
 & Present Bias & Immediate costs and benefits given disproportionate weight over future ones. & 5 & 17 \\
\addlinespace[2pt]\multicolumn{5}{@{}l}{\textbf{Social \& Authority} ($n_{\mathrm{leaves}}=24$)} \\\nopagebreak
\noalign{\needspace{5\baselineskip}}
\multirow{3}{0.17\textwidth}{Authority Deference And Compliance} & Authority Bias & Instructions from high-status figures followed with reduced critical scrutiny. & 3 & 24 \\
 & Obedience To Authority & Harmful directives complied with when issued by a legitimate authority. & 3 & 24 \\
 & Psychological Reactance & Perceived threat to freedom triggers motivational opposition to compliance. & 3 & 24 \\
\noalign{\needspace{8\baselineskip}}
\multirow{6}{0.17\textwidth}{Conformity And Consensus} & Bandwagon Effect & Beliefs or behaviors adopted because others appear to hold them. & 6 & 24 \\
 & Conformity Bias & Individual judgment suppressed in favor of group consensus position. & 6 & 24 \\
 & False Consensus Effect & Own views overestimated as typical of the general population. & 6 & 24 \\
 & Groupthink & Cohesive groups suppress dissent in favor of consensus and harmony. & 6 & 24 \\
 & Pluralistic Ignorance & Private disagreement masked by mistaken belief that others agree. & 6 & 24 \\
 & Social Proof & Others' behavior treated as informative evidence about correct action. & 6 & 24 \\
\noalign{\needspace{9\baselineskip}}
\multirow{7}{0.17\textwidth}{Group Identity And Intergroup Judgment} & Hostile Media Effect & Balanced coverage perceived as biased against one's own group. & 7 & 24 \\
 & Identity Protective Cognition & Reasoning serves to protect group-consistent identities over accuracy. & 7 & 24 \\
 & In Group Bias & Members of one's own group evaluated more favorably than outsiders. & 7 & 24 \\
 & In Group Favoritism & Resources and trust preferentially allocated to in-group members. & 7 & 24 \\
 & Out Group Homogeneity Bias & Out-group members perceived as more similar to each other than in-group. & 7 & 24 \\
 & Reactive Devaluation & Proposals devalued because they originate from an adversarial source. & 7 & 24 \\
 & Stereotype Bias & Group membership used as a heuristic to infer individual attributes. & 7 & 24 \\
\noalign{\needspace{6\baselineskip}}
\multirow{4}{0.17\textwidth}{Liking Similarity Reciprocity And Scarcity} & Liking Bias & Greater compliance granted to those who are liked or admired. & 4 & 24 \\
 & Reciprocity Norm & Obligation to return favors drives compliance with subsequent requests. & 4 & 24 \\
 & Scarcity Effect & Limited availability increases perceived value and urgency of acquisition. & 4 & 24 \\
 & Similarity Attraction Effect & Perceived similarity inflates trust and positive evaluation of a source. & 4 & 24 \\
\noalign{\needspace{6\baselineskip}}
\multirow{4}{0.17\textwidth}{Responsibility Distribution And Public Self} & Bystander Effect & Helping behavior suppressed in groups due to diffused responsibility. & 4 & 24 \\
 & Diffusion Of Responsibility & Individual accountability reduced when others are present to act. & 4 & 24 \\
 & Evaluation Apprehension & Performance altered by concern about being observed and judged. & 4 & 24 \\
 & Social Desirability Bias & Responses skewed toward socially acceptable over truthful answers. & 4 & 24 \\
\addlinespace[2pt]\multicolumn{5}{@{}l}{\textbf{Self-Assessment} ($n_{\mathrm{leaves}}=22$)} \\
\noalign{\needspace{7\baselineskip}}
\multirow{5}{0.17\textwidth}{Attribution Biases} & Actor Observer Bias & Own behavior attributed to situation; others' behavior attributed to disposition. & 5 & 22 \\
 & Correspondence Bias & Behavior attributed to stable traits despite known situational constraints. & 5 & 22 \\
 & Defensive Attribution & Blame assigned to victims to preserve belief in a just world. & 5 & 22 \\
 & Fundamental Attribution Error & Situational causes underweighted relative to dispositional ones in explaining others. & 5 & 22 \\
 & Self Serving Bias & Successes attributed internally; failures attributed to external circumstances. & 5 & 22 \\
\noalign{\needspace{8\baselineskip}}
\multirow{6}{0.17\textwidth}{Confidence Calibration Errors} & Dunning Kruger Effect & Low competence associated with inflated self-assessment of ability. & 6 & 22 \\
 & Overclaiming Bias & Knowledge claimed in domains where none actually exists. & 6 & 22 \\
 & Overconfidence Effect & Confidence in judgments systematically exceeds their accuracy. & 6 & 22 \\
 & Overestimation & Absolute performance overestimated relative to objective measurement. & 6 & 22 \\
 & Overplacement & Own performance ranked above that of peers beyond what accuracy warrants. & 6 & 22 \\
 & Overprecision & Confidence intervals on uncertain quantities are too narrow. & 6 & 22 \\
\noalign{\needspace{6\baselineskip}}
\multirow{4}{0.17\textwidth}{Explanatory Forecasting And Projection Illusions} & Affective Forecasting Error & Future emotional reactions predicted with systematic inaccuracy. & 4 & 22 \\
 & Illusion Of Explanatory Depth & Causal understanding of complex systems overestimated until probed. & 4 & 22 \\
 & Planning Fallacy & Project completion time and cost underestimated despite past overruns. & 4 & 22 \\
 & Projection Bias & Current preferences assumed to persist unchanged into the future. & 4 & 22 \\
\noalign{\needspace{4\baselineskip}}
\multirow{2}{0.17\textwidth}{Moral Self Regulation Distortions} & Moral Licensing & Prior virtuous acts permit subsequent ethical transgressions without guilt. & 2 & 22 \\
 & Self Licensing Effect & Good deeds in one domain license less-virtuous behavior in another. & 2 & 22 \\
\noalign{\needspace{7\baselineskip}}
\multirow{5}{0.17\textwidth}{Reflective Blind Spots And Closure} & Bias Blind Spot & Own susceptibility to cognitive biases underestimated relative to others. & 5 & 22 \\
 & Cognitive Miserliness & System-2 reasoning avoided in favor of low-effort heuristic processing. & 5 & 22 \\
 & Naive Realism & Own perceptions assumed to be objective while others' are seen as biased. & 5 & 22 \\
 & Need For Closure & Preference for definite answers drives premature epistemic termination. & 5 & 22 \\
 & Outcome Bias & Quality of decision judged by its outcome rather than the process. & 5 & 22 \\
\addlinespace[2pt]\multicolumn{5}{@{}l}{\textbf{Interface \& Automation} ($n_{\mathrm{leaves}}=14$)} \\
\noalign{\needspace{7\baselineskip}}
\multirow{5}{0.17\textwidth}{Automation And Algorithm Reliance} & Algorithm Appreciation & Algorithmic recommendations weighted more heavily than human judgment. & 5 & 14 \\
 & Algorithm Aversion & Algorithmic advice discounted after observing any imperfection. & 5 & 14 \\
 & Automation Bias & Erroneous automated outputs accepted without human verification. & 5 & 14 \\
 & Automation Complacency & Reduced vigilance due to over-reliance on automated monitoring systems. & 5 & 14 \\
 & Machine Heuristic & Machines assumed to be objective and free from human-like errors. & 5 & 14 \\
\noalign{\needspace{8\baselineskip}}
\multirow{6}{0.17\textwidth}{Option Structure And Order Effects} & Choice Overload & Excessive options reduce satisfaction and decision quality. & 6 & 14 \\
 & Compromise Effect & Middle option preferred when bracketed by two more extreme alternatives. & 6 & 14 \\
 & Decoy Effect & Dominated option shifts preference toward the dominating target option. & 6 & 14 \\
 & Partition Dependence & Judged frequencies influenced by how categories are partitioned. & 6 & 14 \\
 & Position Bias & Item position in a list influences selection independent of content. & 6 & 14 \\
 & Serial Position Effect & Items at the beginning and end of a list are better remembered. & 6 & 14 \\
\noalign{\needspace{5\baselineskip}}
\multirow{3}{0.17\textwidth}{Routine Interaction And Alert Response} & Capture Error & Familiar routine captures control from an intended non-routine action. & 3 & 14 \\
 & Habituation & Repeated stimulus exposure reduces attentional and behavioral response. & 3 & 14 \\
 & Post Completion Error & Task goal achieved leads to forgetting a final required procedural step. & 3 & 14 \\
\addlinespace[2pt]\multicolumn{5}{@{}l}{\textbf{Privacy \& Disclosure} ($n_{\mathrm{leaves}}=6$)} \\
\noalign{\needspace{4\baselineskip}}
\multirow{2}{0.17\textwidth}{Privacy And Disclosure Tradeoffs} & Privacy Paradox & Stated privacy concerns contradicted by actual disclosure behavior. & 2 & 6 \\
 & Third Person Effect & Others seen as more vulnerable to persuasion than oneself. & 2 & 6 \\
\noalign{\needspace{6\baselineskip}}
\multirow{4}{0.17\textwidth}{Self Presentation Visibility And Audience Misperception} & False Uniqueness Effect & Own views or abilities perceived as more distinctive than they actually are. & 4 & 6 \\
 & Illusion Of Transparency & Internal states assumed to be more visible to others than they are. & 4 & 6 \\
 & Online Disinhibition Effect & Reduced accountability online lowers thresholds for disclosure and risk. & 4 & 6 \\
 & Spotlight Effect & Amount of attention others pay to oneself overestimated. & 4 & 6 \\
\addlinespace[2pt]\multicolumn{5}{@{}l}{\textbf{Affect \& Mood} ($n_{\mathrm{leaves}}=7$)} \\
\noalign{\needspace{4\baselineskip}}
\multirow{2}{0.17\textwidth}{Hot Cold Gaps And State Projection} & Hot Cold Empathy Gap & Current visceral state limits accurate prediction of different-state preferences. & 2 & 7 \\
 & Visceral Bias & Immediate drive state dominates decision-making over deliberative reasoning. & 2 & 7 \\
\noalign{\needspace{4\baselineskip}}
\multirow{2}{0.17\textwidth}{Incidental Affect And Valuation} & Affect Heuristic & Current mood used as information when making evaluative judgments. & 2 & 7 \\
 & Misattribution Of Affect & Incidental emotional arousal attributed to an unrelated target object. & 2 & 7 \\
\noalign{\needspace{5\baselineskip}}
\multirow{3}{0.17\textwidth}{Mood Congruence And State Dependence} & Mood Congruent Judgment & Evaluations biased toward congruence with current affective state. & 3 & 7 \\
 & Mood Congruent Memory & Memory retrieval facilitated for information matching current mood. & 3 & 7 \\
 & State Dependent Memory & Information encoded in one state recalled more easily in the same state. & 3 & 7 \\
\end{longtable}
